\chapter{\MakeUppercase{Future Scope}}
\setcounter{equation}{0}

\section{{Authors' Future Directions}}

The selected paper investigates the prediction of stroke patients' responses to Brain--Computer Interface (BCI)-based rehabilitation using multimodal patient data and Explainable Artificial Intelligence (XAI). Any specific future-work directions attributed to the authors should be verified against the original paper and cited before final submission.

\section{{Possible Improvements Proposed in This Seminar}}

\begin{itemize}
    \item \textbf{Larger and Diverse Datasets:} Evaluate the models using larger patient populations from different clinical settings to improve generalizability.

    \item \textbf{Long-Term Outcome Prediction:} Investigate whether baseline patient characteristics can predict rehabilitation outcomes over longer follow-up periods.

    \item \textbf{Additional Data Sources:} Incorporate neuroimaging, electrophysiological measurements, and other relevant clinical features to investigate potential improvements in prediction.

    \item \textbf{BCI Control Ability:} Include factors affecting patients' ability to control the BCI system to investigate whether they improve treatment-response prediction.

    \item \textbf{Model Validation and Optimization:} Evaluate the models using independent datasets and investigate methods to improve predictive performance and reliability.

    \item \textbf{Enhanced Explainability:} Explore additional XAI techniques to improve understanding of model predictions and their relevance to rehabilitation planning.
\end{itemize}

These are possible research extensions proposed in this seminar. They should not be interpreted as confirmed future-work statements by the original authors or as improvements already demonstrated by the study.